\documentclass[12pt,a4paper,oneside]{report}

% --------------------------------------------------------------- Packages
\usepackage[utf8]{inputenc}
\usepackage[T1]{fontenc}
\usepackage[english]{babel}
\usepackage{lmodern}
\usepackage[margin=1in]{geometry}
\usepackage{setspace}
\usepackage{graphicx}
\usepackage{caption}
\usepackage{subcaption}
\usepackage{booktabs}
\usepackage{longtable}
\usepackage{array}
\usepackage{tabularx}
\usepackage{amsmath}
\usepackage{amssymb}
\usepackage{xcolor}
\usepackage{titlesec}
\usepackage{tocloft}
\usepackage{fancyhdr}
\usepackage{enumitem}
\usepackage[hidelinks]{hyperref}

% --------------------------------------------------------------- Settings
\onehalfspacing
\graphicspath{{figures/}}
\setlength{\parskip}{6pt}
\setlength{\parindent}{0pt}

\pagestyle{fancy}
\fancyhf{}
\fancyhead[L]{\small Brain Tumor Detection}
\fancyhead[R]{\small \thepage}
\renewcommand{\headrulewidth}{0.4pt}

\titleformat{\chapter}[display]
  {\normalfont\huge\bfseries}{\chaptertitlename\ \thechapter}{20pt}{\Huge}

\captionsetup{font=small,labelfont=bf,justification=centering}

\usepackage{float}

\begin{document}

\begin{titlepage}
  \centering
  \vspace*{2cm}
  {\LARGE\bfseries Brain Tumor Detection from MRI Images \\[6pt]
   using an Ensemble of CNN, Transfer Learning \\[6pt]
   and Swin Transformer Models\par}
  \vspace{2cm}
  {\large 8th Semester Project Report\par}
  \vspace{0.6cm}
  {\itshape submitted to\par}
  \vspace{0.4cm}
  {\Large\bfseries <UNIVERSITY\_NAME>\par}
  \vspace{1.0cm}
  In partial fulfilment of the requirements \\
  for the award of the degree of \\[8pt]
  {\large\bfseries Bachelor of Technology\par}
  \textit{in} \\[6pt]
  {\large\bfseries <SPECIALIZATION>\par}
  \vspace{1.0cm}
  \textit{by} \\[6pt]
  <STUDENT\_NAME\_1> (Roll No. <ROLL\_1>) \\[4pt]
  <STUDENT\_NAME\_2> (Roll No. <ROLL\_2>) \\[4pt]
  <STUDENT\_NAME\_3> (Roll No. <ROLL\_3>) \\[4pt]
  <STUDENT\_NAME\_4> (Roll No. <ROLL\_4>) \\
  \vspace{1.0cm}
  \textit{Under the guidance of} \\[6pt]
  {\bfseries <SUPERVISOR\_NAME>, <SUPERVISOR\_DESIGNATION>\par}
  \vspace{1.0cm}
  Department of <DEPARTMENT\_NAME> \\[6pt]
  {\bfseries <UNIVERSITY\_NAME>\par}
  \vspace{0.5cm}
  <SUBMISSION\_MONTH\_YEAR>
\end{titlepage}
\pagenumbering{roman}

\chapter*{Declaration of Authorship}
\addcontentsline{toc}{chapter}{Declaration of Authorship}

We hereby declare that the project report entitled \textit{"Brain Tumor
Detection from MRI Images using an Ensemble of CNN, Transfer Learning and
Swin Transformer Models"} is an authentic record of our own work carried
out at the Department of <DEPARTMENT\_NAME>, <UNIVERSITY\_NAME>, during
the 8th semester of the academic year <ACADEMIC\_YEAR> under the
supervision of <SUPERVISOR\_NAME>, <SUPERVISOR\_DESIGNATION>.

We further declare that the matter embodied in this project report has not
been submitted by us for the award of any other degree or diploma of this
or any other Institute or University. All the information has been
obtained and presented in accordance with academic rules and ethical
conduct. We also declare that, as required by these rules and conduct, we
have fully cited and referenced all materials and results that are not
original to this work. The dataset used in this project is publicly
available under its original license; no patient-identifiable data was
collected, generated, or stored.

\vspace{0.8cm}
Place: <UNIVERSITY\_ADDRESS> \\
Date: <SUBMISSION\_MONTH\_YEAR>

\vspace{0.8cm}
\textbf{Signatures:}

\vspace{1.2cm}
\noindent
\begin{tabularx}{\textwidth}{XX}
\hrulefill & \hrulefill \\
<STUDENT\_NAME\_1> & <STUDENT\_NAME\_2> \\
Roll No: <ROLL\_1> & Roll No: <ROLL\_2> \\[1.2cm]
\hrulefill & \hrulefill \\
<STUDENT\_NAME\_3> & <STUDENT\_NAME\_4> \\
Roll No: <ROLL\_3> & Roll No: <ROLL\_4> \\
\end{tabularx}

\chapter*{Certificate of Recommendation}
\addcontentsline{toc}{chapter}{Certificate of Recommendation}

This is to certify that the Dissertation Report entitled \textit{"Brain
Tumor Detection from MRI Images using an Ensemble of CNN, Transfer
Learning and Swin Transformer Models"} submitted by <STUDENT\_NAME\_1>,
<STUDENT\_NAME\_2>, <STUDENT\_NAME\_3>, and <STUDENT\_NAME\_4> to
<UNIVERSITY\_NAME> is a record of bonafide project work carried out by
them under my supervision and guidance, and is worthy of consideration for
the award of the degree of Bachelor of Technology (B.Tech) in
<SPECIALIZATION>.

\vspace{1.5cm}
\noindent
\textbf{<SUPERVISOR\_NAME>,} \\
Project Supervisor, \\
<SUPERVISOR\_DESIGNATION>, Department of <DEPARTMENT\_NAME>, \\
<UNIVERSITY\_NAME>.

\vspace{1.0cm}
\textbf{Approved By:}

\vspace{1.0cm}
\noindent
\textbf{<HOD\_NAME>,} \\
HoD / Reviewer, Department of <DEPARTMENT\_NAME>, \\
<UNIVERSITY\_NAME>.

\chapter*{Acknowledgement}
\addcontentsline{toc}{chapter}{Acknowledgement}

We would like to first express our sincere gratitude to our project
supervisor, <SUPERVISOR\_NAME>, <SUPERVISOR\_DESIGNATION>, Department of
<DEPARTMENT\_NAME>, <UNIVERSITY\_NAME>. Their continued guidance,
technical insight, and constructive feedback during every stage of this
project were instrumental in shaping the final outcome. Their willingness
to discuss difficult engineering decisions, including the diagnosis of
subtle numerical-stability issues during mixed-precision training, helped
us approach the work with both rigour and curiosity.

We take this opportunity to express our gratitude to all faculty members
of the Department of <DEPARTMENT\_NAME> for their support, the
encouragement they provided during our coursework, and the foundation
they laid in machine learning, computer vision, and software engineering.

We acknowledge the open-source community whose tools made this work
feasible: the PyTorch and timm library maintainers, the authors of
pytorch-grad-cam, the Streamlit and FastAPI teams, and Masoud Nickparvar
for releasing the Brain Tumor MRI Dataset on Kaggle under a permissive
license that supports academic research.

Finally, we thank our parents and families for their unceasing
encouragement, patience, and support throughout the semester.

\chapter*{Abstract}
\addcontentsline{toc}{chapter}{Abstract}

The accurate and timely classification of brain tumors from magnetic
resonance imaging (MRI) is a clinically important problem, yet manual
interpretation by radiologists is time-consuming, expertise-intensive,
and subject to inter-observer variability. Reported inter-rater agreement
on tumor type lies in the 70 to 85 percent range, while access to expert
neuroradiologists remains uneven across geographies. There is a clear
case for AI-assisted decision support that is fast, consistent, and
explainable.

This project presents an end-to-end deep-learning system that classifies
T1-weighted brain MRI slices into four clinical categories: glioma,
meningioma, pituitary tumor, and no tumor. The system uses an ensemble of
three architecturally diverse models. The first model is a custom
convolutional neural network (CNN) of approximately 1.2 million
parameters, trained from random initialisation. The second is
EfficientNet-B0, adapted via two-stage transfer learning from ImageNet,
totalling about 4.0 million parameters. The third is Swin-Tiny, a
hierarchical vision transformer of 27.5 million parameters that uses
windowed self-attention and is similarly fine-tuned in two stages. The
three models' softmax probability outputs are combined by a
weighted-average ensemble whose mixing weights were obtained by
exhaustive grid search on a held-out validation set.

The system was developed and evaluated on the publicly available Brain
Tumor MRI Dataset (Nickparvar, Kaggle, 7,200 images across four classes).
A stratified 80/20 split partitions the 5,600 training images into 4,480
training and 1,120 validation samples, while 1,600 testing images remain
held out and are evaluated exactly once. On this held-out test set, the
strongest individual model (EfficientNet-B0) achieves 94.94\% accuracy
and a macro-averaged ROC-AUC of 0.991; the three-model ensemble achieves
94.69\% accuracy at 0.991 AUC. We report an honest empirical finding:
that ensembling matches but does not exceed the dominant base learner,
and we discuss it through the lens of the dominant-model problem in
ensemble learning.

Beyond accuracy, the system provides per-prediction Grad-CAM heatmaps
for all three models, a low-confidence threshold gate that explicitly
recommends radiologist review for uncertain cases, and four deployment
channels: a Streamlit web interface, a FastAPI REST backend, ONNX
exports that yield a 2 to 3.5x CPU speedup, and a Docker container.

\vspace{0.5cm}
\textbf{Keywords:} Brain Tumor Classification, Magnetic Resonance Imaging,
Convolutional Neural Networks, Transfer Learning, Swin Transformer,
Ensemble Learning, Grad-CAM, Computer-Aided Diagnosis.

\tableofcontents
\listoffigures
\listoftables
\clearpage
\pagenumbering{arabic}

\chapter{Introduction}

\section{Medical Imaging and Computer-Aided Diagnosis}

Brain tumors are abnormal growths of cells that form inside the cranium. They can originate from brain tissue itself (primary tumors) or spread to the brain from other parts of the body (secondary or metastatic tumors). The World Health Organization (WHO) classifies primary brain tumors into over 150 subtypes based on cell of origin, molecular markers, and grade of malignancy. Among these, gliomas, meningiomas, and pituitary adenomas are the three most commonly encountered categories in clinical practice.

Magnetic Resonance Imaging (MRI) is the preferred diagnostic modality for brain tumors because it provides excellent soft-tissue contrast without exposing the patient to ionising radiation. A standard brain MRI examination produces dozens of two-dimensional slices across multiple sequences (T1-weighted, T2-weighted, FLAIR, and contrast-enhanced T1), each of which must be visually inspected by a neuroradiologist. Manual interpretation of a complete scan typically takes 15 to 30 minutes per patient, and reported inter-observer agreement on tumor type ranges from 70 to 85 percent, meaning that even experienced radiologists disagree in a non-trivial fraction of cases~\cite{bouhafra2025}.

Computer-Aided Diagnosis (CAD) systems aim to assist radiologists by providing automated second-reader opinions. Over the past decade, deep learning, and convolutional neural networks (CNNs) in particular, have transformed the CAD landscape. A well-trained deep learning model can analyse an MRI slice in milliseconds, provide a confidence-calibrated prediction, and highlight the regions of the image that most influenced its decision~\cite{ranjbarzadeh2025}.

India records approximately 28,000 new brain tumor cases per year according to the International Association of Cancer Registries and the Indian Journal of Neurology (2025), with some estimates reaching 40,000 to 50,000 cases annually. Brain and central nervous system tumors account for roughly 2 to 3 percent of all cancers in the country. Meanwhile, India has fewer than one radiologist per 100,000 people, with the shortage being far more severe in rural and semi-urban areas~\cite{iarc2025}.

\begin{figure}[H]
  \centering
  \includegraphics[width=0.95\textwidth]{Figure_08_system_flow_diagram.png}
  \caption{High-level system flow: input MRI through preprocessing, three parallel models, weighted ensemble, to predicted class with confidence and Grad-CAM heatmaps.}
  \label{fig:system-flow}
\end{figure}

\section{Objectives of the Project}

\subsection{Primary Objective}

To develop a deep-learning-based brain tumor classification system that takes a single T1-weighted brain MRI slice as input and predicts one of four diagnostic categories (glioma, meningioma, pituitary tumor, or no tumor) with high accuracy, calibrated confidence, and visual explainability, using an ensemble of three architecturally diverse neural networks.

\subsection{Secondary Objectives}

\begin{itemize}[leftmargin=*]
  \item Objective 1: To design and train a custom CNN from scratch to serve as a baseline model and to provide architectural diversity for the ensemble.
  \item Objective 2: To adapt a pretrained EfficientNet-B0 model for brain tumor classification using a two-stage transfer learning strategy that maximises accuracy while avoiding catastrophic forgetting of pretrained features.
  \item Objective 3: To fine-tune a Swin-Tiny vision transformer for the same classification task, introducing self-attention-based architectural diversity alongside the CNN-based models.
  \item Objective 4: To combine the three models into a weighted-average ensemble whose weights are determined by exhaustive grid search on the validation set.
  \item Objective 5: To implement Gradient-weighted Class Activation Mapping (Grad-CAM) for all three models to provide visual explanations.
  \item Objective 6: To build a Streamlit web interface and a FastAPI REST backend that wrap the trained models and expose real-time prediction with Grad-CAM overlays.
  \item Objective 7: To export all three models to ONNX format for cross-platform deployment.
  \item Objective 8: To rigorously evaluate the system on a held-out test set of 1,600 images and report honest metrics including accuracy, macro F1-score, macro ROC-AUC, per-class confusion matrices, and inference latency.
\end{itemize}

\section{Challenges in Brain Tumor Classification}

\begin{itemize}[leftmargin=*]
  \item \textbf{High intra-class variability:} Tumors of the same type can appear very different across patients depending on size, location, growth stage, and scanner settings.
  \item \textbf{Inter-class visual similarity:} Gliomas and meningiomas can appear remarkably similar on a single T1-weighted slice.
  \item \textbf{Limited training data:} Medical imaging datasets are inherently small compared to the millions of images available for natural image tasks like ImageNet.
  \item \textbf{Scanner and protocol heterogeneity:} Images acquired on different MRI scanners produce systematically different intensity profiles.
  \item \textbf{Explainability and clinical trust:} Clinicians are unlikely to adopt a diagnostic tool that functions as a black box.
  \item \textbf{Calibrated uncertainty:} A system that outputs overconfident incorrect predictions is worse than one that acknowledges uncertainty.
\end{itemize}

\section{Methodologies}

The methodology of this project follows a structured, multi-stage pipeline that proceeds from data acquisition and preparation through model training, ensemble construction, explainability, and deployment.

\begin{itemize}[leftmargin=*]
  \item \textbf{Data Acquisition:} The Brain Tumor MRI Dataset (Nickparvar, Kaggle) is used, containing 7,200 labelled MRI images across four classes (see Chapter 3).
  \item \textbf{Preprocessing and Augmentation:} All images are converted to RGB, resized to 224 by 224 pixels, and normalised using ImageNet statistics. Details are in Chapter 3.
  \item \textbf{Model Training:} Three architecturally diverse models are trained (Chapter 5).
  \item \textbf{Ensemble Construction:} Soft voting, weighted average, and stacking are evaluated. Optimal weights are found by exhaustive grid search (Chapter 5).
  \item \textbf{Explainability:} Grad-CAM heatmaps are generated for each base model (Chapter 5).
  \item \textbf{Deployment:} Streamlit web app, FastAPI REST service, ONNX exports, and a Docker container (Chapters 4 and 5).
\end{itemize}

\chapter{Literature Survey}

\section{Overview of the Field}

Two recent systematic reviews provide a useful starting point. Bouhafra and El Bahi (2025)~\cite{bouhafra2025} reviewed 60 papers on deep-learning brain tumor classification published between 2020 and January 2024, covering transfer learning, autoencoders, transformers, and attention mechanisms. Ranjbarzadeh et al. (2025)~\cite{ranjbarzadeh2025} published a comprehensive survey of explainable AI and vision-transformer approaches specifically for brain tumor detection, examining performance, datasets, preprocessing, and interpretability across the 2020 to 2024 window. These reviews establish three dominant research themes: (1) transfer learning from ImageNet-pretrained backbones (predominantly EfficientNet and ResNet families), (2) vision-transformer architectures (especially the Swin family), and (3) ensemble methods that combine multiple architectures.

Section~2.2 reviews ten of the most directly relevant papers in detail. Each subsection presents the paper's methodology, dataset, headline results, and the specific way in which the paper informed the design of our own project. Section~2.3 then consolidates the reviewed work in a single comparison table.

\section{Review of Related Work}

\subsection{Babu Vimala et al. (2023): Hybrid Deep Learning with EfficientNet Variants}

Babu Vimala et al.~\cite{babuvimala2023} proposed a transfer-learning approach using fine-tuned EfficientNet variants (B0 through B4) on the CE-MRI Figshare dataset for three-class brain tumor classification (glioma, meningioma, pituitary). Their methodology employed a two-step refinement: initialising models with ImageNet weights, then adding custom top layers for tumor classification. Grad-CAM visualisation was incorporated to highlight attention maps. The best-performing model, EfficientNetB2, achieved 99.06\% test accuracy with 98.79\% F1-score. This study established EfficientNet as a strong baseline for medical image classification when combined with appropriate fine-tuning and data augmentation. Our project draws on this work in its use of two-stage transfer learning for the EfficientNet-B0 base model.

\subsection{Islam et al. (2024): BrainNet Optimised EfficientNet Architecture}

Islam et al.~\cite{islam2024} presented BrainNet, an optimised EfficientNet pipeline evaluated on the CE-MRI dataset of 3,064 T1-weighted images. The study compared all EfficientNet variants from B0 to B7 and incorporated advanced preprocessing and augmentation techniques. The best model, EfficientNetB3, achieved 99.69\% test accuracy. The paper emphasises the role of learning-rate scheduling and progressive unfreezing in achieving top performance on small medical imaging datasets. This work directly informed our project's two-stage training schedule, where Stage 1 trains only the head and Stage 2 unfreezes the backbone with a 10x smaller learning rate.

\subsection{Alnowami et al. (2024): EfficientNetv2 with Attention Mechanisms}

Alnowami et al.~\cite{alnowami2024eff} enhanced the EfficientNetv2 architecture with a Global Attention Mechanism (GAM) and Efficient Channel Attention (ECA) for MRI-based brain tumor classification. The integration of attention mechanisms allowed the model to focus on salient features within complex MRI images, improving classification accuracy on a four-class dataset. The model achieved 99.76\% accuracy, setting a new benchmark at the time of publication. The work also incorporated Grad-CAM visualisation for clinical interpretability. This study demonstrates that attention-augmented CNNs can rival full transformer architectures at lower computational cost, an insight relevant to our choice of EfficientNet-B0 over heavier transformer models given the 4 GB VRAM constraint.

\subsection{Liu et al. (2021): Original Swin Transformer}

Liu et al.~\cite{liu2021swin} introduced the Swin Transformer, a hierarchical vision transformer that uses shifted windows for efficient self-attention. Unlike earlier vision transformers that compute attention over the entire image (quadratic cost), Swin partitions the image into non-overlapping windows and computes attention only within each window. In alternating layers, the windows are shifted by half their size, allowing information to flow across window boundaries. This design yields linear computational complexity in image size while preserving global modelling capability. The Swin-Tiny variant achieves 81.3\% top-1 accuracy on ImageNet at only 28 million parameters. This paper is the foundation for our third base model.

\subsection{Haq et al. (2024): Modified Swin Transformer with HSW-MSA and ResMLP}

Haq et al.~\cite{haq2024swin} proposed a modified Swin Transformer for brain tumor diagnosis, introducing a Hybrid Shifted Windows Multi-Head Self-Attention (HSW-MSA) module along with a Residual MLP (ResMLP) replacing the traditional MLP layer. These modifications aimed to improve classification accuracy, reduce memory usage, and simplify training complexity. The proposed model was evaluated on a publicly available four-class brain MRI dataset using transfer learning and data augmentation. The model achieved a remarkable 99.92\% accuracy, surpassing previous CNN-based methods. This paper demonstrates the strong potential of transformer architectures for medical image classification and motivated our inclusion of Swin-Tiny as one of the three base models in our ensemble.

\subsection{Alsubai et al. (2024): Hybrid Swin Transformer and ResNet50V2}

Alsubai et al.~\cite{alsubai2024swin} proposed a hybrid architecture combining the Swin Transformer with ResNet50V2 for enhanced brain tumor classification from MRI images. The Swin Transformer captures local-and-global hierarchical features through windowed self-attention, while ResNet50V2 provides complementary deep convolutional features. The two streams are fused through transfer learning to leverage pretrained weights from both backbones. The hybrid architecture demonstrates that combining local feature extraction (CNN) with global context modelling (transformer) yields robust classification across diverse MRI acquisition conditions. This dual-stream philosophy parallels our project's ensemble approach, which combines three architecturally distinct models.

\subsection{Khan et al. (2023): Weighted Average Ensemble Deep Learning}

Khan et al.~\cite{khan2023ensemble} proposed a weighted average ensemble model combining three deep-learning feature spaces: a VGG19 model (transfer learning), a standard CNN, and a CNN with data augmentation. The optimal combination of weights was found by exhaustive grid search. The dataset used was The Cancer Genome Atlas (TCGA) lower-grade glioma collection with 3,929 MRI images. The ensemble outperformed each individual model in accuracy, precision, and F1-score, demonstrating that combining models with different feature spaces can compensate for individual model weaknesses. This paper directly motivates our project's grid-search approach to ensemble-weight optimisation.

\subsection{Abdella et al. (2025): Majority Voting Ensemble of Seven CNNs}

Abdella et al.~\cite{abdella2025majority} evaluated seven pretrained CNN architectures (including GoogLeNet and Inception-v3) for four-class brain tumor classification on T1-weighted MRI images. Each model was trained using two optimisers (SGDM and Adam) and evaluated on a public dataset split into training (70\%), validation (10\%), and testing (20\%) subsets. A majority voting ensemble aggregating predictions from all 14 trained models achieved 99.8\% accuracy with AUC values above 0.997 for all tumor classes, outperforming every standalone model. The work highlights that model diversity is the key ingredient for successful ensembling, a finding that aligns with our project's deliberate architectural diversity (CNN + EfficientNet + Swin) rather than combining variants of the same family.

\subsection{Ali et al. (2024): Grad-CAM with ResNet50 for Clinical Explainability}

Ali et al.~\cite{ali2024gradcam} addressed the critical challenge of interpretability in deep-learning brain tumor detection by integrating Grad-CAM with ResNet50. The methodology trained ResNet50 on an MRI dataset with data augmentation, then applied Grad-CAM to generate heatmap visualisations highlighting regions of the MRI most influential to the model's predictions. The model achieved 98.52\% testing accuracy with precision-recall metrics exceeding 98\%. The work demonstrates that combining high accuracy with visual explanations significantly increases clinical trust and aids radiologists in verifying that the model focuses on relevant pathological markers. Our project adopts the same Grad-CAM approach, extended to all three base models and the predicted class of the ensemble.

\subsection{Naser and Deen (2023): Multi-CAM Comparative Study}

Naser and Deen~\cite{naser2023gradcam} developed a lightweight Explainable Deep Learning framework comparing three explainability techniques, Class Activation Mapping (CAM), Grad-CAM, and Grad-CAM++, applied to three models: pretrained VGG-19, scratch VGG-19, and EfficientNet. Evaluation on two benchmark MRI datasets (multi-class and binary) showed that pretrained VGG-19 with Grad-CAM outperformed other combinations both in classification accuracy and in the visual quality of localisation heatmaps. This study is one of the few that quantitatively compares different CAM variants for medical imaging, and confirms that Grad-CAM (which we use in our project) consistently produces the most clinically meaningful heatmaps.

\section{Summary of Reviewed Approaches}

Table~\ref{tab:litreview-summary} consolidates the ten reviewed papers, showing for each one the method, dataset, headline accuracy, and the primary way in which the paper informed our project's design.

\begin{table}[H]
  \centering
  \caption{Summary comparison of the reviewed literature.}
  \label{tab:litreview-summary}
  \begin{tabular}{p{2.4cm}cp{4.0cm}cp{4.2cm}}
    \toprule
    \textbf{Reference} & \textbf{Year} & \textbf{Method} & \textbf{Accuracy} & \textbf{Contribution to our project} \\
    \midrule
    Babu Vimala et al. & 2023 & EfficientNetB0-B4 + Grad-CAM & 99.06\% & Two-step transfer learning baseline \\
    Islam et al. & 2024 & EfficientNetB0-B7 with augmentation & 99.69\% & Progressive unfreezing schedule \\
    Alnowami et al. & 2024 & EfficientNetv2 + GAM + ECA & 99.76\% & Attention can rival transformers \\
    Liu et al. & 2021 & Swin Transformer (foundational) & 81.3\%\textsuperscript{$\dagger$} & Shifted-window self-attention \\
    Haq et al. & 2024 & Modified Swin + HSW-MSA + ResMLP & 99.92\% & Transformer-on-MRI feasibility \\
    Alsubai et al. & 2024 & Swin + ResNet50V2 hybrid & high & CNN + transformer fusion idea \\
    Khan et al. & 2023 & Weighted ensemble VGG19 + CNN & best of 3 & Grid-search ensemble weights \\
    Abdella et al. & 2025 & Majority voting (7 CNNs) & 99.8\% & Diversity drives ensembles \\
    Ali et al. & 2024 & ResNet50 + Grad-CAM & 98.52\% & Clinical-trust framework \\
    Naser \& Deen & 2023 & VGG-19 + CAM/Grad-CAM/++ & high & Grad-CAM beats CAM and Grad-CAM++ \\
    \bottomrule
  \end{tabular}
\end{table}

\textsuperscript{$\dagger$} ImageNet top-1 accuracy; the Swin Transformer is the foundational architecture, not a brain-MRI-specific result.

\section{Market Research}

The global burden of brain tumors is substantial and growing. According to the Global Cancer Observatory (GLOBOCAN 2022), approximately 308,000 new cases of brain and central nervous system tumors were diagnosed worldwide in 2020, with an age-standardised global incidence rate of 5.57 per 100,000 population~\cite{globocan2022}.

In India, the picture is particularly acute. The International Association of Cancer Registries and the Indian Journal of Neurology (2025) report 28,000 new brain tumor cases per year, though some expert estimates range as high as 40,000 to 50,000 cases annually. Brain tumors are the second most common cancer site in children aged 0 to 14 years~\cite{iarc2025}.

Against this disease burden, India faces a severe shortage of diagnostic specialists. The country has fewer than one radiologist per 100,000 people, with the deficit being most acute in tier-2 and tier-3 cities and in rural regions. In this context, AI-assisted diagnostic tools that can provide reliable second-reader opinions in real time are not a luxury but a clinical necessity~\cite{iarc2025}.

\section{Research Gap and Limitations of Existing Systems}

Despite the substantial body of work, several limitations persist in the existing literature:

\begin{itemize}[leftmargin=*]
  \item \textbf{Overstated accuracy:} Many studies report only validation-set accuracy without a truly held-out test set.
  \item \textbf{Limited architectural diversity in ensembles:} Most ensemble studies combine variants of the same architecture family. Few works deliberately combine architecturally distinct model families.
  \item \textbf{Insufficient explainability:} While Grad-CAM is mentioned in several papers, few provide per-model heatmap comparisons or analyse failure modes.
  \item \textbf{Lack of deployment engineering:} Most published works stop at model training and metric reporting.
  \item \textbf{Absence of honest negative results:} The scenario where a dominant base learner makes ensembling counterproductive is rarely acknowledged.
\end{itemize}

\section{Motivation for Project Selection}

The selection of this project was motivated by the convergence of a pressing clinical need, a well-characterised public dataset, and a desire to address the research gaps identified above. The project seeks to answer:

\begin{itemize}[leftmargin=*]
  \item Can an ensemble of three architecturally diverse models yield better classification than any single model alone?
  \item How large is the gap between validation-set and test-set accuracy?
  \item Can Grad-CAM heatmaps provide clinically meaningful visual explanations, and do different model architectures attend to different anatomical regions?
  \item Can the entire system be packaged as a deployable web application suitable for clinical demonstration?
\end{itemize}

\chapter{Dataset and Data Preparation}

\section{Overview}

A well-characterised and reproducible dataset is the foundation of any supervised deep-learning system. This chapter describes the dataset source, the clinical meaning of each class, the rationale for selecting this dataset, the exploratory analysis performed, the train/validation/test split, the preprocessing pipeline, and the augmentation strategy.

\section{Dataset Source and Acquisition}

The dataset used throughout this project is the publicly available Brain Tumor MRI Dataset, released by Masoud Nickparvar on the Kaggle platform~\cite{nickparvar2022}. The dataset is freely accessible at:

\url{https://www.kaggle.com/datasets/masoudnickparvar/brain-tumor-mri-dataset}

The collection is a curated aggregation of three earlier public brain MRI resources: the figshare brain tumor dataset, the SARTAJ brain tumor dataset, and the Br35H brain tumor detection dataset, re-organised into a single unified four-class taxonomy with a consistent Training/Testing folder layout. After per-class balancing, the dataset contains a total of 7,200 images split as 5,600 in the Training folder and 1,600 in the Testing folder.

\section{Class Categories and Clinical Meaning}

\textbf{Glioma:} Gliomas are neoplasms that arise from the glial cells of the central nervous system. They constitute the most common form of malignant primary brain tumor in adults. The WHO grades gliomas on a scale from I (least aggressive) to IV (most aggressive). On T1-weighted MRI, gliomas frequently appear as ill-defined intra-axial masses with surrounding vasogenic oedema.

\textbf{Meningioma:} Meningiomas arise from the arachnoid cap cells of the meninges. They are the most common primary brain tumor overall and are typically benign and slow-growing. On MRI, they appear as well-circumscribed extra-axial masses attached to the dura. Their appearance on a single T1-weighted axial slice can be visually similar to that of a glioma.

\textbf{Pituitary tumor:} Pituitary adenomas are growths arising from the pituitary gland at the base of the brain. The overwhelming majority are benign. On MRI, pituitary lesions appear within or just superior to the sella turcica, at the very base of the brain.

\textbf{No tumor:} This class comprises MRI slices with no detectable mass lesion. Its inclusion is essential for clinical usability: a diagnostic system must be capable of confidently identifying a normal scan.

\begin{figure}[H]
  \centering
  \includegraphics[width=0.95\textwidth]{Figure_02_sample_mri_per_class.jpg}
  \caption{Representative MRI slices, one per class.}
  \label{fig:samples}
\end{figure}

\section{Rationale for Dataset Selection}

\begin{itemize}[leftmargin=*]
  \item \textbf{Four-class taxonomy aligned with the clinical problem.}
  \item \textbf{Class balance.} After per-class augmentation, the Training folder contains a near-uniform number of images per class.
  \item \textbf{Sufficient volume for transfer learning.} With approximately 7,200 images, the dataset is large enough to fine-tune ImageNet-pretrained backbones.
  \item \textbf{Public and permissive licensing.} The dataset contains no patient identifiers.
  \item \textbf{Established benchmark.} The dataset has been used in dozens of peer-reviewed studies since 2021.
  \item \textbf{Heterogeneity.} The aggregation of three upstream collections yields a corpus spanning multiple scanners and intensity profiles.
\end{itemize}

\section{Dataset Statistics}

The Training folder is partitioned via stratified random sampling into model-training and validation subsets at an 80/20 ratio. The Testing folder is retained as a strictly held-out evaluation set.

\begin{table}[H]
  \centering
  \caption{Per-class image counts in the Training, Validation, and Test splits.}
  \label{tab:split-counts}
  \begin{tabular}{lccccc}
    \toprule
    \textbf{Split} & \textbf{glioma} & \textbf{meningioma} & \textbf{notumor} & \textbf{pituitary} & \textbf{Total} \\
    \midrule
    Training (full) & 1,400 & 1,400 & 1,400 & 1,400 & 5,600 \\
    ~~Train (80\%) & 1,120 & 1,120 & 1,120 & 1,120 & 4,480 \\
    ~~Validation (20\%) & 280 & 280 & 280 & 280 & 1,120 \\
    Test (held out) & 400 & 400 & 400 & 400 & 1,600 \\
    \textbf{Grand total} & \textbf{1,800} & \textbf{1,800} & \textbf{1,800} & \textbf{1,800} & \textbf{7,200} \\
    \bottomrule
  \end{tabular}
\end{table}

\begin{figure}[H]
  \centering
  \includegraphics[width=0.95\textwidth]{Figure_01_class_distribution.png}
  \caption{Per-class image counts across the Training, Validation, and Test splits.}
  \label{fig:class-dist}
\end{figure}

\section{Exploratory Data Analysis}

\textbf{File integrity:} Every image file was opened with the Python Imaging Library. No corrupt files were encountered.

\textbf{Image modes:} Of the 7,200 images, 2,585 are stored as RGB, 1,892 as single-channel grayscale, and 3 as four-channel RGBA. A uniform conversion to RGB is applied during preprocessing.

\textbf{Pixel dimensions:} Image dimensions vary widely, from 150 by 168 pixels to 1,375 by 1,446 pixels. All images are resized to 224 by 224 during preprocessing.

\textbf{Pixel-intensity distribution:} Per-class pixel-intensity histograms overlap substantially across all four classes, confirming that no class is separable by simple thresholding.

\begin{figure}[H]
  \centering
  \includegraphics[width=0.95\textwidth]{Figure_03_image_size_distribution.png}
  \caption{Distribution of original image dimensions.}
  \label{fig:image-sizes}
\end{figure}

\begin{figure}[H]
  \centering
  \includegraphics[width=0.95\textwidth]{Figure_05_pixel_intensity_histograms.png}
  \caption{Per-class pixel-intensity histograms showing substantial overlap.}
  \label{fig:intensity-hist}
\end{figure}

\section{Train / Validation / Test Split Strategy}

\begin{itemize}[leftmargin=*]
  \item \textbf{Training set (4,480 images, 62.2\%):} Used for gradient updates.
  \item \textbf{Validation set (1,120 images, 15.6\%):} Used for early stopping, hyperparameter selection, and ensemble-weight grid search.
  \item \textbf{Test set (1,600 images, 22.2\%):} Held out entirely until all hyperparameters are frozen, then evaluated exactly once.
\end{itemize}

The split is performed by \texttt{sklearn.model\_selection.train\_test\_split} with \texttt{stratify=labels} and \texttt{random\_state=42}, guaranteeing bit-exact reproducibility.

\section{Image Preprocessing Pipeline}

\begin{table}[H]
  \centering
  \caption{Image preprocessing transforms applied to all data.}
  \label{tab:preprocessing}
  \begin{tabular}{cll}
    \toprule
    \textbf{Step} & \textbf{Transform} & \textbf{Parameters} \\
    \midrule
    1 & Convert to RGB & PIL.Image.convert('RGB') \\
    2 & Resize & 224 x 224 px, bilinear \\
    3 & To tensor & float32, range [0, 1] \\
    4 & Normalise & ImageNet mean / std \\
    \bottomrule
  \end{tabular}
\end{table}

\begin{figure}[H]
  \centering
  \includegraphics[width=0.95\textwidth]{Figure_06_preprocessing_effect.jpg}
  \caption{Visual effect of the deterministic preprocessing pipeline.}
  \label{fig:preproc-effect}
\end{figure}

\section{Data Augmentation Strategy}

Augmentation is applied only to training images and only at training time.

\begin{table}[H]
  \centering
  \caption{Training-time augmentation policy.}
  \label{tab:augmentation}
  \begin{tabular}{lll}
    \toprule
    \textbf{Augmentation} & \textbf{Probability / Range} & \textbf{Clinical Justification} \\
    \midrule
    Random horizontal flip & p = 0.5 & Brain is bilaterally symmetric \\
    Random rotation & +/- 15 deg & Simulates patient head tilt \\
    Random affine (translate) & +/- 5\% & Field-of-view variation \\
    Random affine (scale) & +/- 5\% & Magnification differences \\
    Brightness jitter & +/- 0.20 & Scanner-gain variation \\
    Contrast jitter & +/- 0.20 & Intensity-window variation \\
    Random erasing & p=0.25, area 2-10\% & Forces global reasoning \\
    \bottomrule
  \end{tabular}
\end{table}

\begin{figure}[H]
  \centering
  \includegraphics[width=0.95\textwidth]{Figure_07_augmentation_examples.jpg}
  \caption{Examples of training-time data augmentation.}
  \label{fig:augment-examples}
\end{figure}

\section{Ethical Considerations and Licensing}

The dataset is publicly distributed for academic and research use and contains no patient-identifiable information. The output of the trained system is intended for academic demonstration only and is not represented as a medical device.

\chapter{Project Planning}

\section{Overview}

Project planning defines the structured roadmap for the development and completion of the Brain Tumor Detection system. The planning process focuses on dividing the project into well-defined phases to ensure systematic progress.

\section{Current Status of the Project}

At the present stage, the project is fully implemented and operational. The following activities have been successfully completed:

\begin{itemize}[leftmargin=*]
  \item Problem domain identification and objectives formulation
  \item Literature survey and analysis of existing systems
  \item Dataset acquisition, EDA, and preprocessing pipeline design
  \item Design and training of three deep learning models
  \item Ensemble construction with grid-search weight optimisation
  \item Rigorous evaluation on the held-out test set
  \item Grad-CAM explainability implementation
  \item Streamlit web interface and FastAPI REST backend
  \item ONNX model export with inference verification
  \item Docker containerisation for cloud deployment
  \item Complete documentation and report preparation
\end{itemize}

\section{System Architecture}

The Brain Tumor Detection system follows a layered architecture designed for modularity, maintainability, and ease of deployment. The system is divided into four logical layers.

\textbf{Presentation Layer:} Provides two user-facing interfaces. The Streamlit web application offers an interactive demo with drag-and-drop image upload, real-time prediction display, confidence metrics, per-model probability breakdowns, and Grad-CAM heatmap overlays. The FastAPI REST backend exposes the same capability as a JSON API.

\textbf{Application Layer:} The \texttt{BrainTumorPredictor} class in \texttt{src/inference.py} is the single entry point for all prediction requests. Both the Streamlit app and the FastAPI backend call this same class.

\textbf{Model Layer:} Contains the three trained model architectures, the ensemble combiner, the Grad-CAM implementation, and the training engine.

\textbf{Data Layer:} The dataset, cached split CSVs, training outputs (plots, confusion matrices, Grad-CAM images), and TensorBoard logs.

\begin{figure}[H]
  \centering
  \includegraphics[width=0.95\textwidth]{Figure_08_system_flow_diagram.png}
  \caption{Layered architecture and data flow of the Brain Tumor Detection platform.}
  \label{fig:arch}
\end{figure}

\section{Identified Challenges and System Limitations}

\begin{itemize}[leftmargin=*]
  \item \textbf{GPU memory constraint:} 4 GB VRAM limit on the NVIDIA GTX 1650.
  \item \textbf{Mixed-precision numerical instability:} fp16 triggered a NaN bug on Turing architecture.
  \item \textbf{Dominant-model effect:} EfficientNet-B0 is significantly stronger than the other two models.
  \item \textbf{Single MRI sequence:} The system classifies from a single T1-weighted slice only.
\end{itemize}

\section{Strategy to Address Identified Challenges}

\begin{itemize}[leftmargin=*]
  \item \textbf{Bfloat16 mixed precision} resolved the NaN issue.
  \item \textbf{Aggressive but clinically valid augmentation} mitigates overfitting.
  \item \textbf{Architectural diversity} maximises uncorrelated errors.
  \item \textbf{Grad-CAM integration} addresses the clinical-trust barrier.
  \item \textbf{Threshold gate} flags uncertain predictions for radiologist review.
\end{itemize}

\section{Implementation Roadmap}

\begin{table}[H]
  \centering
  \caption{Phases of the entire project work.}
  \label{tab:phases}
  \begin{tabular}{cll}
    \toprule
    \textbf{Phase} & \textbf{Module} & \textbf{Tasks} \\
    \midrule
    1 & Data + EDA & Download dataset, integrity check, EDA, splits \\
    2 & Custom CNN & Design 4-block VGG-style CNN, train from scratch \\
    3 & EfficientNet-B0 & Two-stage transfer learning \\
    4 & Swin-Tiny & Two-stage fine-tuning with very small Stage 2 LR \\
    5 & Ensemble & Grid search over weights, select best \\
    6 & Evaluation & Test set: accuracy, F1, AUC, CMs, ROC \\
    7 & Explainability & Grad-CAM for all three architectures \\
    8 & Deployment & Streamlit, FastAPI, ONNX, Docker \\
    9 & Documentation & Report, viva guide, architecture diagrams \\
    \bottomrule
  \end{tabular}
\end{table}

\section{Project Scheduling and Gantt Chart}

The project was executed over a single academic semester. Phases 2 through 4 (model training) were partially parallelised. The ensemble and evaluation phases were strictly sequential, as they depend on the availability of all three trained models.

\begin{figure}[H]
  \centering
  \includegraphics[width=0.95\textwidth]{Figure_31_gantt_chart.png}
  \caption{Project schedule across the semester.}
  \label{fig:gantt}
\end{figure}

\chapter{Project Description}

\section{Software Model}

The Brain Tumor Detection system is built using a modular Python package architecture. The \texttt{src/} directory contains all production code, organised into single-responsibility modules. The \texttt{app/} directory contains the two user-facing interfaces. The \texttt{notebooks/} directory contains seven Jupyter notebooks that reproduce the entire experimental pipeline.

\begin{table}[H]
  \centering
  \caption{Key technology choices and justifications.}
  \label{tab:tech-choices}
  \begin{tabular}{lll}
    \toprule
    \textbf{Component} & \textbf{Choice} & \textbf{Reason} \\
    \midrule
    Framework & PyTorch 2.12 & Industry standard for research \\
    Mixed precision & bfloat16 & Avoids fp16 cuDNN NaN bug \\
    Optimiser & AdamW & Strong baseline for small datasets \\
    Loss function & CrossEntropyLoss & Classes are balanced \\
    Batch size & 32 & Fits Swin-Tiny in 4 GB VRAM with AMP \\
    LR scheduler & CosineAnnealingLR & Smooth convergence \\
    Random seed & 42 (all RNGs) & Bit-exact reproducibility \\
    \bottomrule
  \end{tabular}
\end{table}

\section{Software Requirements Specification (SRS)}

\subsection{Introduction}

The Brain Tumor Detection system is a deep-learning-based diagnostic support tool that classifies brain MRI slices into four categories.

\subsection{General Description}

The system is a client-server application with a Python backend and a web-based frontend. Both interfaces (Streamlit and FastAPI) wrap the same \texttt{BrainTumorPredictor} class.

\subsection{Functional Requirements}

The system shall:

\begin{itemize}[leftmargin=*]
  \item Accept a brain MRI image in JPG, PNG, or BMP format.
  \item Preprocess the input image (RGB convert, resize, normalise).
  \item Run the preprocessed image through all three trained models.
  \item Combine the three probability vectors using the saved ensemble weights.
  \item Return the predicted class and confidence score.
  \item Optionally generate Grad-CAM heatmap overlays for each base model.
  \item Flag predictions below a user-configurable threshold (default 0.85) for radiologist review.
  \item Display results in the Streamlit UI including: predicted class, confidence, probability bar chart, per-model probability table, and three Grad-CAM heatmaps.
  \item Expose a \texttt{/predict} REST endpoint returning JSON.
\end{itemize}

\subsection{Interface Requirements}

\begin{itemize}[leftmargin=*]
  \item \textbf{User Interface:} Web-based, accessible via any modern browser; drag-and-drop image upload; sidebar with model info.
  \item \textbf{Hardware Interface:} Runs on standard desktop or laptop with or without a dedicated GPU.
  \item \textbf{Software Interface:} Python 3.11+, PyTorch 2.12, timm, Streamlit, FastAPI, ONNX Runtime.
\end{itemize}

\subsection{Performance Requirements}

\begin{itemize}[leftmargin=*]
  \item Single-image inference shall complete within 300 ms on GPU (including Grad-CAM) or 1 second on CPU.
  \item Model loading shall complete within 10 seconds on first startup.
  \item The Streamlit interface shall remain responsive during prediction.
\end{itemize}

\subsection{Design Constraints}

\begin{itemize}[leftmargin=*]
  \item Open-source frameworks only.
  \item All model weights stored locally; no external API calls.
  \item Deployable on standard institutional infrastructure.
\end{itemize}

\subsection{Non-Functional Attributes}

\begin{itemize}[leftmargin=*]
  \item \textbf{Reliability:} Deterministic predictions for the same input.
  \item \textbf{Security:} No user data stored or transmitted externally.
  \item \textbf{Maintainability:} Modular package structure.
  \item \textbf{Scalability:} FastAPI supports concurrent requests; ONNX enables edge deployment.
\end{itemize}

\section{Functional Specification}

\subsection{Deep Learning Models}

\textbf{Model 1: Custom CNN.}

A VGG-style architecture consisting of four convolutional blocks. Each block contains two stacked 3x3 convolutions, each followed by Batch Normalisation and ReLU activation, and a 2x2 max-pooling layer. The channel count doubles at each block: 32, 64, 128, 256. After the fourth block, the feature map (256 x 14 x 14) passes through Global Average Pooling, then a two-layer fully-connected classifier with dropout. Weight initialisation uses Kaiming normal. Total parameters: 1,206,628.

\begin{figure}[H]
  \centering
  \includegraphics[width=0.95\textwidth]{Figure_09_architecture_custom_cnn.png}
  \caption{Custom CNN architecture.}
  \label{fig:arch-cnn}
\end{figure}

\textbf{Model 2: EfficientNet-B0 Transfer Learning.}

EfficientNet-B0 was chosen over ResNet-50, EfficientNet-B3, DenseNet-121, and ConvNeXt-Tiny on the basis of its accuracy-per-parameter efficiency and 4 GB VRAM compatibility. The pretrained backbone is loaded from torchvision; its 1000-class classifier head is replaced with Dropout(0.3) and Linear(1280 to 4). Training proceeds in two stages: Stage 1 (4 epochs, LR = 0.001) with the backbone frozen, then Stage 2 (8 epochs, LR = 0.0001) with the backbone unfrozen. Total parameters: 4,012,672.

\begin{figure}[H]
  \centering
  \includegraphics[width=0.95\textwidth]{Figure_10_architecture_efficientnet_b0.png}
  \caption{EfficientNet-B0 transfer learning architecture with two-stage training schedule.}
  \label{fig:arch-eff}
\end{figure}

\textbf{Model 3: Swin-Tiny Transformer.}

Swin-Tiny is a hierarchical vision transformer that splits the input image into 4x4 patches and processes them through four stages of windowed self-attention blocks. The shifted-window mechanism allows information to flow across window boundaries. Two-stage training mirrors the EfficientNet recipe but with a Stage 2 LR of 0.00002 (50x smaller) because transformer weights are extremely sensitive. Total parameters: 27,522,430.

\begin{figure}[H]
  \centering
  \includegraphics[width=0.95\textwidth]{Figure_11_architecture_swin_tiny.png}
  \caption{Swin-Tiny architecture with patch embedding and four hierarchical stages.}
  \label{fig:arch-swin}
\end{figure}

\begin{table}[H]
  \centering
  \caption{Architectural summary of the three base models.}
  \label{tab:arch-summary}
  \begin{tabular}{llll}
    \toprule
    \textbf{Property} & \textbf{Custom CNN} & \textbf{EfficientNet-B0} & \textbf{Swin-Tiny} \\
    \midrule
    Total params & 1.21 M & 4.01 M & 27.52 M \\
    Architecture & VGG conv blocks & Depthwise-sep + SE & Shifted-window MSA \\
    Pretrained & No (scratch) & Yes (ImageNet) & Yes (ImageNet) \\
    Stage 1 LR & 0.001 & 0.001 & 0.001 \\
    Stage 2 LR & N/A & 0.0001 & 0.00002 \\
    Stage 1 epochs & 20 & 4 & 3 \\
    Stage 2 epochs & N/A & 8 & 4 \\
    \bottomrule
  \end{tabular}
\end{table}

\subsection{Mathematical Model and Loss Functions}

All three models are trained with the standard cross-entropy loss function. For a single sample with true class $y$ and predicted probability $p_y$ for that class, the loss is:

$$ \mathcal{L} = -\log(p_y) $$

where $p_y$ is the softmax probability assigned to the correct class. Since the four classes are perfectly balanced, no class-weighting is applied to the loss.

The optimiser used is AdamW with $\beta_1 = 0.9$, $\beta_2 = 0.999$, weight decay $= 10^{-4}$. The learning rate follows CosineAnnealingLR, decaying smoothly from the initial value to $\eta_{\min} = \eta_0 \times 10^{-2}$.

\subsection{Ensemble Decision Logic}

Each trained model outputs a softmax probability vector $\mathbf{p}_i \in \mathbb{R}^4$. The weighted average ensemble computes:

$$ \mathbf{p}_{\text{ens}} = w_{\text{cnn}} \mathbf{p}_{\text{cnn}} + w_{\text{tl}} \mathbf{p}_{\text{tl}} + w_{\text{swin}} \mathbf{p}_{\text{swin}} $$

with $w_i \ge 0$ and $\sum_i w_i = 1$. Optimal weights were found by exhaustive grid search over the simplex at step 0.05 (231 candidates), yielding $[w_{\text{cnn}}, w_{\text{tl}}, w_{\text{swin}}] = [0.0, 0.9, 0.1]$.

\begin{figure}[H]
  \centering
  \includegraphics[width=0.95\textwidth]{Figure_27_ensemble_weight_heatmap.png}
  \caption{Validation accuracy across the ensemble-weight grid search.}
  \label{fig:ens-heatmap}
\end{figure}

\subsection{Grad-CAM Explainability Generation}

Gradient-weighted Class Activation Mapping (Grad-CAM) produces a coarse localisation heatmap highlighting the regions of an input image most important for a specific predicted class~\cite{selvaraju2017gradcam}. For the custom CNN and EfficientNet, the target layer is the last convolutional layer before global average pooling. For Swin-Tiny, the target is the LayerNorm at the end of the final transformer stage, with a reshape operation to handle the channels-last activation format.

\begin{figure}[H]
  \centering
  \includegraphics[width=0.95\textwidth]{Figure_28_gradcam_correct.jpg}
  \caption{Representative Grad-CAM overlays for correctly classified samples (one per class).}
  \label{fig:gradcam-correct}
\end{figure}

\section{Design Specification}

\subsection{Use-Case Diagrams}

The system supports three primary actor roles.

\begin{figure}[H]
  \centering
  \includegraphics[width=0.95\textwidth]{Figure_15_usecase_radiologist.png}
  \caption{Use-case diagram: Radiologist / Clinician role.}
  \label{fig:uc-rad}
\end{figure}

\begin{figure}[H]
  \centering
  \includegraphics[width=0.95\textwidth]{Figure_16_usecase_researcher.png}
  \caption{Use-case diagram: Researcher / Developer role.}
  \label{fig:uc-res}
\end{figure}

\begin{figure}[H]
  \centering
  \includegraphics[width=0.95\textwidth]{Figure_17_usecase_admin.png}
  \caption{Use-case diagram: Administrator role.}
  \label{fig:uc-adm}
\end{figure}

\subsection{Data Flow Diagrams}

\begin{figure}[H]
  \centering
  \includegraphics[width=0.95\textwidth]{Figure_18_dfd_training.png}
  \caption{Data Flow Diagram for the model training pipeline.}
  \label{fig:dfd-train}
\end{figure}

\begin{figure}[H]
  \centering
  \includegraphics[width=0.95\textwidth]{Figure_30_dfd_inference.png}
  \caption{Data Flow Diagram for single-image inference.}
  \label{fig:dfd-inf}
\end{figure}

\subsection{Data Dictionary}

\begin{table}[H]
  \centering
  \caption{Ensemble Configuration Schema.}
  \label{tab:dd-ensemble}
  \begin{tabular}{lll}
    \toprule
    \textbf{Field} & \textbf{Type} & \textbf{Description} \\
    \midrule
    method & string & weighted, soft\_voting, or stacking \\
    weights & list of float & Mixing weights (sum to 1.0) \\
    val\_accuracy & float & Validation accuracy \\
    models & list of string & ['cnn', 'transfer', 'swin'] \\
    \bottomrule
  \end{tabular}
\end{table}

\begin{table}[H]
  \centering
  \caption{Prediction Result Schema.}
  \label{tab:dd-pred}
  \begin{tabular}{lll}
    \toprule
    \textbf{Field} & \textbf{Type} & \textbf{Description} \\
    \midrule
    predicted\_class & string & glioma / meningioma / pituitary / notumor \\
    confidence & float & Ensemble top-class probability \\
    ensemble\_probs & dict & Class name to probability \\
    model\_probs & dict of dict & Per-model breakdown \\
    inference\_time\_ms & float & Wall-clock inference time \\
    gradcams & dict or None & Heatmap arrays (if requested) \\
    \bottomrule
  \end{tabular}
\end{table}

\chapter{Implementation Issues}

\section{Dataset Variability and Image Quality}

The most immediate challenge was the heterogeneity of the source images. The Nickparvar dataset aggregates three upstream collections, so images vary in colour mode, resolution, intensity profile, and contrast. While this heterogeneity benefits model robustness, it required careful preprocessing to ensure that every image was in a consistent format. The three RGBA images required explicit handling.

\section{Model Selection and Hyperparameter Sensitivity}

Five candidate transfer-learning backbones were evaluated: ResNet-50, EfficientNet-B0, EfficientNet-B3, DenseNet-121, and ConvNeXt-Tiny. EfficientNet-B0 was selected for its superior accuracy-per-parameter ratio and 4 GB VRAM compatibility.

The Swin Transformer was extremely sensitive to the Stage 2 learning rate. A Stage 2 LR of 0.0001 (the EfficientNet value) caused the transformer's validation accuracy to collapse within three epochs. Reducing to 0.00002 (50x smaller) resolved the issue.

\begin{figure}[H]
  \centering
  \includegraphics[width=0.95\textwidth]{Figure_12_cnn_training_curves.png}
  \caption{Custom CNN training and validation curves.}
  \label{fig:tc-cnn}
\end{figure}

\begin{figure}[H]
  \centering
  \includegraphics[width=0.95\textwidth]{Figure_13_efficientnet_training_curves.png}
  \caption{EfficientNet-B0 training curves across both stages.}
  \label{fig:tc-eff}
\end{figure}

\begin{figure}[H]
  \centering
  \includegraphics[width=0.95\textwidth]{Figure_14_swin_training_curves.png}
  \caption{Swin-Tiny training curves across both stages.}
  \label{fig:tc-swin}
\end{figure}

\section{GPU Memory Constraints and Mixed Precision}

The entire training pipeline was constrained to 4 GB of VRAM on an NVIDIA GTX 1650 (Turing architecture, compute capability 7.5). Mixed-precision training using fp16 was initially employed but produced NaN values in the loss after approximately 5 epochs during Swin-Tiny training.

Root-cause analysis identified a cuDNN numerical instability in LayerNorm operations when computed in fp16 on Turing GPUs. The solution was to switch to bfloat16, which maintains the same dynamic range as float32 (8 exponent bits) while using only 16 bits total~\cite{micikevicius2018mixed}.

\section{Class Confusion and Generalisation Gap}

Analysis of confusion matrices reveals that the primary source of classification error is the glioma-meningioma confusion pair. The custom CNN misclassifies 17\% of gliomas as meningiomas; EfficientNet-B0 reduces this to 1\%; Swin-Tiny to 6\%. This confusion is clinically expected: on a single T1-weighted axial slice, these two tumor types can appear visually similar.

\begin{figure}[H]
  \centering
  \includegraphics[width=0.95\textwidth]{Figure_19_confusion_matrix_cnn.png}
  \caption{Confusion matrix on the test set: Custom CNN.}
  \label{fig:cm-cnn}
\end{figure}

\begin{figure}[H]
  \centering
  \includegraphics[width=0.95\textwidth]{Figure_20_confusion_matrix_efficientnet.png}
  \caption{Confusion matrix on the test set: EfficientNet-B0.}
  \label{fig:cm-eff}
\end{figure}

\begin{figure}[H]
  \centering
  \includegraphics[width=0.95\textwidth]{Figure_21_confusion_matrix_swin.png}
  \caption{Confusion matrix on the test set: Swin-Tiny.}
  \label{fig:cm-swin}
\end{figure}

\begin{figure}[H]
  \centering
  \includegraphics[width=0.95\textwidth]{Figure_22_confusion_matrix_ensemble.png}
  \caption{Confusion matrix on the test set: Weighted ensemble.}
  \label{fig:cm-ens}
\end{figure}

\begin{figure}[H]
  \centering
  \includegraphics[width=0.95\textwidth]{Figure_24_per_class_f1_bars.png}
  \caption{Per-class F1 scores across all models.}
  \label{fig:f1-bars}
\end{figure}

\section{Clinical Adoption and Explainability Trust}

From a deployment perspective, the most significant barrier to clinical adoption is not accuracy but trust. The system addresses this through Grad-CAM heatmaps, per-model probability breakdowns, and a configurable confidence threshold.

Qualitative analysis of Grad-CAM outputs on correctly classified samples confirms that all three models attend to anatomically meaningful regions. EfficientNet's heatmaps are the most spatially focused; Swin's are the most diffuse. On misclassified samples, Grad-CAM reveals that the model often attends to the correct anatomical region but misinterprets the visual pattern.

\begin{figure}[H]
  \centering
  \includegraphics[width=0.95\textwidth]{Figure_29_gradcam_misclassified.jpg}
  \caption{Grad-CAM overlays for misclassified samples, illustrating failure modes.}
  \label{fig:gradcam-mis}
\end{figure}

\section{The Dominant-Model Problem in Ensembling}

The most important and unexpected empirical finding is that the three-model ensemble (94.69\% test accuracy) does not outperform the best individual model (EfficientNet-B0, 94.94\% test accuracy). This result contradicts the common expectation that ensembles always improve over individual models.

The result is explained by two factors. First, EfficientNet-B0 already operates near the data's noise ceiling. Second, the custom CNN and Swin-Tiny make many of the same errors as EfficientNet, so averaging their predictions does not cancel error but instead injects noise from the weaker models.

This phenomenon is documented in the ensemble learning literature as the dominant-model problem: variance reduction from ensembling only materialises if base learners have uncorrelated errors~\cite{khan2023ensemble}.

\begin{figure}[H]
  \centering
  \includegraphics[width=0.95\textwidth]{Figure_23_roc_curves_test.png}
  \caption{Receiver Operating Characteristic curves on the test set for all three models and the ensemble.}
  \label{fig:roc}
\end{figure}

\begin{figure}[H]
  \centering
  \includegraphics[width=0.95\textwidth]{Figure_25_final_comparison_bars.png}
  \caption{Final comparison bar chart of test-set accuracy, F1, and macro AUC across models and ensemble.}
  \label{fig:final-bars}
\end{figure}

\begin{figure}[H]
  \centering
  \includegraphics[width=0.95\textwidth]{Figure_26_efficiency_pareto.png}
  \caption{Accuracy versus inference latency trade-off.}
  \label{fig:pareto}
\end{figure}

\begin{table}[H]
  \centering
  \caption{Per-model test-set performance.}
  \label{tab:perf}
  \begin{tabular}{lrrrrr}
    \toprule
    \textbf{Model} & \textbf{Params} & \textbf{Val acc} & \textbf{Test acc} & \textbf{Macro F1} & \textbf{Macro AUC} \\
    \midrule
    Custom CNN & 1.21 M & 0.9036 & 0.8275 & 0.8221 & 0.9513 \\
    EfficientNet-B0 & 4.01 M & 0.9893 & \textbf{0.9494} & \textbf{0.9482} & \textbf{0.9908} \\
    Swin-Tiny & 27.52 M & 0.9670 & 0.9194 & 0.9171 & 0.9865 \\
    Ensemble (weighted) & 32.74 M & 0.9893 & 0.9469 & 0.9456 & 0.9907 \\
    \bottomrule
  \end{tabular}
\end{table}

\begin{table}[H]
  \centering
  \caption{Ensemble grid-search results (top five weight combinations).}
  \label{tab:grid}
  \begin{tabular}{ccccc}
    \toprule
    \textbf{Rank} & \textbf{$w_{\text{cnn}}$} & \textbf{$w_{\text{tl}}$} & \textbf{$w_{\text{swin}}$} & \textbf{Val accuracy} \\
    \midrule
    1 & 0.00 & 0.90 & 0.10 & 0.9893 \\
    2 & 0.00 & 0.95 & 0.05 & 0.9893 \\
    3 & 0.00 & 1.00 & 0.00 & 0.9893 \\
    4 & 0.05 & 0.90 & 0.05 & 0.9884 \\
    5 & 0.00 & 0.85 & 0.15 & 0.9875 \\
    \bottomrule
  \end{tabular}
\end{table}

\chapter{Conclusion, Summary and Future Scope}

\section{Summary of the Work}

This project designed, implemented, and rigorously evaluated a complete deep-learning system for four-class brain tumor classification from T1-weighted MRI images. The system combines a custom CNN (1.2 M params), an EfficientNet-B0 with two-stage transfer learning (4.0 M params), and a Swin-Tiny vision transformer (27.5 M params), combined through a weighted-average ensemble with weights $[0.0, 0.9, 0.1]$ found by grid search.

The system was trained on the Brain Tumor MRI Dataset (Nickparvar, Kaggle), with a stratified 80/20 split producing 4,480 training and 1,120 validation images, and 1,600 test images held out and evaluated exactly once.

Beyond model training, the project delivered four deployment interfaces (Streamlit, FastAPI, ONNX, Docker), Grad-CAM explainability, a configurable confidence threshold gate, and comprehensive documentation.

\section{Conclusion}

The dominant individual model, EfficientNet-B0, achieves \textbf{94.94\% accuracy} on the held-out test set with a macro-averaged ROC-AUC of \textbf{0.991}. The three-model ensemble achieves 94.69\% accuracy at 0.991 AUC, matching but not exceeding the best individual model. This result is consistent with the dominant-model problem documented in the ensemble learning literature.

Grad-CAM analysis confirms that all three models attend to anatomically meaningful regions. Remaining errors map to explainable failure modes, primarily the glioma-meningioma boundary.

The system is fully deployable through four channels, demonstrating that a complete production-quality medical imaging AI system can be built on modest hardware within an academic semester, provided that careful attention is paid to engineering fundamentals.

\section{Future Scope}

\begin{itemize}[leftmargin=*]
  \item \textbf{3D Volumetric Classification:} Extend the pipeline to process full 3D MRI volumes.
  \item \textbf{Multi-Sequence Fusion:} Incorporate multiple MRI sequences (T1, T2, FLAIR, contrast-enhanced T1).
  \item \textbf{Self-Supervised Pretraining:} Train a foundation model on a large corpus of unlabelled brain MRI data.
  \item \textbf{Test-Time Augmentation:} Average predictions over multiple augmented versions of the same test image.
  \item \textbf{Cross-Institution Validation:} Validate on multi-centre clinical data.
  \item \textbf{Federated Learning:} Train across multiple hospital sites without centralising patient data.
  \item \textbf{LLM-Assisted Report Generation:} Integrate a large language model to generate preliminary radiology reports.
\end{itemize}

\begin{thebibliography}{99}

\bibitem{bouhafra2025}
S. Bouhafra and H. El Bahi,
\textit{"Deep Learning Approaches for Brain Tumor Detection and
Classification Using MRI Images (2020 to 2024): A Systematic Review,"}
Journal of Digital Imaging Informatics in Medicine,
vol. 38, pp. 1403--1433, 2025.

\bibitem{ranjbarzadeh2025}
R. Ranjbarzadeh et al.,
\textit{"Explainable AI and Vision Transformers for Detection and
Classification of Brain Tumor: A Comprehensive Survey,"}
Artificial Intelligence Review, Springer, 2025.

\bibitem{babuvimala2023}
B. Babu Vimala, S. Srinivasan, S. K. Mathivanan et al.,
\textit{"Detection and Classification of Brain Tumor Using Hybrid Deep
Learning Models,"}
Scientific Reports, vol. 13, art. 23029, 2023.

\bibitem{islam2024}
M. T. Islam et al.,
\textit{"BrainNet: Precision Brain Tumor Classification with Optimized
EfficientNet Architecture,"}
International Journal of Intelligent Systems, Wiley, 2024.

\bibitem{alnowami2024eff}
M. Alnowami et al.,
\textit{"Enhancing EfficientNetv2 with Global and Efficient Channel
Attention Mechanisms for Accurate MRI-Based Brain Tumor Classification,"}
Cluster Computing, Springer, 2024.

\bibitem{haq2024swin}
A. U. Haq et al.,
\textit{"A Novel Swin Transformer Approach Utilizing Residual Multi-Layer
Perceptron for Diagnosing Brain Tumors in MRI Images,"}
International Journal of Machine Learning and Cybernetics,
vol. 15, pp. 3579--3597, 2024.

\bibitem{alsubai2024swin}
S. Alsubai et al.,
\textit{"Enhanced Magnetic Resonance Imaging-Based Brain Tumor
Classification with a Hybrid Swin Transformer and ResNet50V2 Model,"}
Applied Sciences, vol. 14, no. 22, art. 10154, 2024.

\bibitem{khan2023ensemble}
M. A. Khan et al.,
\textit{"Weighted Average Ensemble Deep Learning Model for Stratification
of Brain Tumor in MRI Images,"}
Diagnostics, vol. 13, no. 7, art. 1320, 2023.

\bibitem{abdella2025majority}
G. M. Abdella et al.,
\textit{"Majority Voting Ensemble of Deep CNNs for Robust MRI-Based Brain
Tumor Classification,"}
Diagnostics, vol. 15, no. 14, art. 1782, 2025.

\bibitem{ali2024gradcam}
M. Ali et al.,
\textit{"Enhancing Brain Tumor Detection in MRI Images through Explainable
AI Using Grad-CAM with ResNet 50,"}
BMC Medical Imaging, 2024.

\bibitem{naser2023gradcam}
M. A. Naser and M. J. Deen,
\textit{"Explainable Deep Learning Approach for Multi-Class Brain MRI
Tumor Classification and Localization Using Gradient-Weighted Class
Activation Mapping,"}
Information, vol. 14, no. 12, art. 642, 2023.

\bibitem{nickparvar2022}
M. Nickparvar,
\textit{"Brain Tumor MRI Dataset,"}
Kaggle, 2022.
\url{https://www.kaggle.com/datasets/masoudnickparvar/brain-tumor-mri-dataset}

\bibitem{selvaraju2017gradcam}
R. R. Selvaraju et al.,
\textit{"Grad-CAM: Visual Explanations from Deep Networks via
Gradient-based Localization,"}
in Proceedings of the IEEE International Conference on Computer Vision
(ICCV), 2017.

\bibitem{tan2019efficientnet}
M. Tan and Q. V. Le,
\textit{"EfficientNet: Rethinking Model Scaling for Convolutional Neural
Networks,"}
in Proceedings of the International Conference on Machine Learning
(ICML), 2019.

\bibitem{liu2021swin}
Z. Liu et al.,
\textit{"Swin Transformer: Hierarchical Vision Transformer Using Shifted
Windows,"}
in Proceedings of the IEEE International Conference on Computer Vision
(ICCV), 2021.

\bibitem{loshchilov2019adamw}
I. Loshchilov and F. Hutter,
\textit{"Decoupled Weight Decay Regularization,"}
in Proceedings of ICLR, 2019.

\bibitem{iarc2025}
International Association of Cancer Registries (IARC) and Indian Journal
of Neurology, 2025;
BW Healthcare World,
\textit{"Battling the Brain Tumor Burden: Challenges and Advancements in
India,"} 2025.

\bibitem{micikevicius2018mixed}
P. Micikevicius et al.,
\textit{"Mixed Precision Training,"}
in Proceedings of ICLR, 2018.

\bibitem{ahmed2025swrd}
S. Ahmed et al.,
\textit{"Advancing Brain Tumor MRI Classification using SwRD: A Parallel
Swin Transformer-ResNet Approach,"}
Virtual Reality and Intelligent Hardware,
vol. 7, no. 5, pp. 501--522, 2025.

\bibitem{globocan2022}
Global Cancer Observatory (GLOBOCAN),
\textit{"Cancer Today: Estimated Number of New Cases in 2020, Worldwide,"}
International Agency for Research on Cancer, World Health Organization,
2022.

\end{thebibliography}

\end{document}
